\section{Conclusion}\label{sec:concl}

The argument of the essay can be stated in a few sentences. A writing system combines a conventional assignment of marks to linguistic units with a perceptual structure among the marks, and the second is not arbitrary even when the first is. Visual substitution alphabets make the difference tangible, because their assignments are interchangeable in principle and their resemblances are not, and the design of such an alphabet is an optimisation whose answer depends on the typeface, on the notion of resemblance used, and on how many letters are changed. A font that separates stored characters from displayed shapes, together with routes between shapes and a schedule over the positions of a book, turns a single text into a family of editions, and the drifting edition, in which the script changes while the text is read, is the member of the family that asks whether a reader can acquire a script receptively through sustained reading.

The analysis of such editions produced results of three kinds. Some are mathematical: the sizes of the spaces of alphabets, the collapse of linear exchange and its removal by a detour, the avoidance of collision by cycles of length three, the conservation of exposure and the lower bound it gives on the length of a drift, the failure of substitution and shaping to commute unless joining classes are preserved, and the gluing criterion for local keys. Some are computational and exactly as strong as their inputs: the comparison of pairings across four typefaces and the simulation of schedules in a fluid model. And some are hypotheses about readers, which the essay leaves as hypotheses and for which it specifies an empirical programme. The pattern across all of them is that the structure of the problem is governed by things one might first overlook: the dependence of resemblance on the typeface, the inevitability of a shared glyph when two letters exchange along a straight path, the dominance of the rarest letter in the length of a drift, and the possibility of a book that is locally coherent at every point and has no global key.

Several lines of work follow. The weights of the pairing problem should be replaced by measured confusion data, and the optimisation redone. The detour marks and keyframe routes should be designed and tested for recognisability. The exposure model should be extended from its fluid approximation to the actual random sequence of letters in a particular text, with the schedule computed from that text. The coherence conditions of Section~\ref{sec:formal} suggest an experiment with and without a visible phase marker for swap-ending drifts. And the formal treatment should be pushed where it is thinnest, in the relation between the sheaf-theoretic model of conventions and measurable properties of readers. Because the source and the builder are public, each of these can be pursued by others on their own texts and compared, and the editions that result will be part of the evidence for or against the proposition with which the essay began: that a publication can be designed as a passage between representational conventions, and that the reader of such a publication can be carried across.
